Multi-LoRA inference pools must place adapter weights, reusable prompt state,
and request demand together. Conversation affinity preserves prefixes but can
replicate each adapter across enough engines to enlarge their working sets and
reduce effective batching. We describe an admission router that combines
measured token work with demand-sized adapter domains. Independent frontends
prefer a bounded set of replicas for each exact adapter version, preserve
observed cache matches, and escape congested placements using a work budget
expressed in engine batches. No cluster-wide coordinator lies on the custom
selection path.

We evaluate real inference on eight H200 GPUs with a hybrid 27B model,
64 saved adapter versions, and controlled general-serving and independent
single- and multi-turn RL-shaped workloads. In two order-reversed 64-version
trials, domains improve throughput over KV affinity by 59.1--62.8\%, reduce
p95 request latency by 52.3--52.7\%, and reduce GPU adapter loads by
85.5--89.1\%. These gains occur despite a 10.1--11.2 percentage-point decrease in KV
hit ratio. Warmer, smaller working sets accompany larger running batches while
requests are queued. RL-shaped replays show smaller and mixed improvements:
rollout time, individual request latency, and replay throughput can move in
different directions; a 1,024-token control regresses in window throughput and
latency. The results establish a substantial fixed-pool serving
benefit, while exposing the limits of local observations, affinity exceptions,
and replay-based claims about training throughput.
